\documentclass[11pt]{article}
\usepackage[margin=1in]{geometry}
\usepackage{amsmath,amssymb,amsfonts,mathtools}
\usepackage{array,booktabs,enumitem,graphicx,xcolor,hyperref}
\usepackage[T1]{fontenc}
\usepackage[utf8]{inputenc}
\title{On the possibility of \$Z\_N\$ exotic supersymmetry in two dimensional
  Conformal Field Theory}
\author{Source-backed arXiv sample}
\date{}
\begin{document}
\maketitle
\section{Extracted Lists}

The list environments below are extracted from arXiv LaTeX source and wrapped for pdfLaTeX validation.

\begin{itemize}
\item
For $P=5$ we have $K=5$, $k=2,3$. This is the result noticed in the paper
that $\SZ_5$ coincides with the $Z_5$ parafermion.
\item
For $P=4$ no solution appears.
\item
For $P=3$ we have $K=6$, $k=2,4$, thus showing that a realization of $\SZ_3$
is contained in the $Z_6$ parafermionic algebra. This is not surprising as
the $Z_6$ model is known to belong to the unitary minimal series of spin 4/3
algebra, namely for $m=4$.
\item
For $P=2$ there are two solutions: $K=6$, $k=3$ says that the $Z_6$ model is
also supersymmetric, (it belongs indeed also to the superconformal minimal
series for $m=6$) while $K=8$, $k=2,6$ shows {\em two} fields of spin 3/2
for the $Z_8$ model.
\item
For $P=1$ the solution $K=8$, $k=4$ completes the result for $P=2$: the two
spin 3/2 curents are associated to a current of spin 2, thus \{$T,\psi_2,
\psi_4,\psi_6$ form in this case two copies of the $N=1$ superconformal
algebra: the $Z_8$ model is doubly $N=1$ supersymmetric. Another solution
appears for $P=1$ when $K=9$, $k=3,6$.
\end{itemize}

\begin{enumerate}
\item
identify the parameters in the algebra. These are $c$ and the independent
structure constants. Moreover, spin of fields will eventually depend on
a set of integers $M_k$ as in eq.(\ref{delta}).
\item
list all the non-zero 4-point functions and compute $R$ for each of them.
Imposing $R\geq 0$ the integers $M_k$ can be selected to a few possible
choices.
\item
For each choice consider the functions $\corr{i}{\ih}{k}{\kh}$. If at least
one of them has $R=0$, discard the choice. Else, identify the $\corr{i}{\ih}
{k}{\kh}$ functions with lower values of $R$ and, after having constrained
as much as possible the coefficients of the polynomials using duality,
try to fix $c$ and/or some $C_{ij}$ through eqs.(\ref{c},\ref{q},\ref{t}).
If different
correlation functions provide unconsistent values of $c$ or $C_{ij}$,
discard the choice.
\item
Otherwise use the values so obtained to fix as much as possible the other
correlation functions and check all the possible duality constraints. If
somewhere some unconsistency appears, discard the choice. If instead the
choice passes all the checks it defines a consistent associative algebra.
If some parameter is left free, the algebra allows for an infinity of
associative realizations, one for each value of the parameter.
\item
redo steps 3 and 4 for all the choices selected by step 2.
\end{enumerate}
\end{document}
